\section*{Bab \ref{ch:b1:second-law-entropy} --- Hukum Kedua: Entropi}

\begin{solution}{exo:b1:second-law-entropy:1}
$\Delta S = mL/T$: peleburan $334000/273 = \qty{1.22}{kJ/K}$; pendidihan $2260000/
373 = \qty{6.06}{kJ/K}$. Penguapan membebaskan molekulnya satu sama lain
sepenuhnya --- jauh lebih banyak ketakteraturan mikroskopik, dan jauh lebih banyak kalor tiap
\omterm{def:b1:kinetic-theory:temperature}{kelvin}.
\end{solution}

\begin{solution}{exo:b1:second-law-entropy:2}
$\Delta S = mc\ln(T_2/T_1) = 4180\ln(353/293) = \qty{0.78}{kJ/K}$ --- sebuah fungsi
keadaan: tak bergantung pada caranya (entropi yang \emph{tercipta} tidak begitu).
\end{solution}

\begin{solution}{exo:b1:second-law-entropy:3}
$\Delta S = nR\ln2 = \qty{5.8}{J/K}$. Isotermal reversibel: $Q = nRT\ln2$
dari rendamannya pada $T$, $S_{\mathrm{tukar}} = nR\ln2$, yang tercipta $0$. Joule:
$Q = 0$, $S_{\mathrm{tukar}} = 0$, yang tercipta \qty{5.8}{J/K}.
\end{solution}

\begin{solution}{exo:b1:second-law-entropy:4}
$S = nC_{V,m}\ln T + nR\ln V$ tetap: $T^{C_{V,m}}V^R$ tetap, yakni
$TV^{R/C_{V,m}} = TV^{\gamma - 1}$ tetap. Pemampatan yang tiba-tiba: $Q = 0$, sehingga
$S_{\mathrm{tukar}} = 0$ dan $\Delta S = S_{\mathrm{cipta}}$; dengan $T$: $300 \to
643$ K, $V$: $24.9 \to 10.7$ L: $\Delta S = 20.8\ln(643/300) + 8.314\ln(10.7/
24.9) = 15.9 - 7.0 = \qty{+8.9}{J/K}$ yang tercipta.
\end{solution}

\begin{solution}{exo:b1:second-law-entropy:5}
Baloknya: $450\ln(293/373) = \qty{-109}{J/K}$; danaunya: $+36000/293 = \qty{+123}{J/K}$;
semestanya: $+\qty{14}{J/K}$.
\end{solution}

\begin{solution}{exo:b1:second-law-entropy:6}
$T_f = \qty{350}{K}$; $S = 4180[\ln(350/400) + \ln(350/300)] = 4180 \times
0.0206 = \qty{86}{J/K}$, semuanya tercipta. $(T_1 + T_2)^2 - 4T_1T_2 = (T_1 - T_2)^2
\geq 0$, sehingga argumen logaritmanya $\geq 1$.
\end{solution}

\begin{solution}{exo:b1:second-law-entropy:7}
$S = Q(1/T_c - 1/T_h)$: $1000(1/280 - 1/300) = \qty{0.24}{J/K}$; dari tungku ke
ruangannya: $1000(1/300 - 1/1000) = \qty{2.3}{J/K}$, sepuluh kali lebih besar: makin besar
selisih suhu yang diseberangi sebuah aliran kalor, makin tak reversibel ia ---
turunan dari tungku ke ruanganlah tempat energi bermutu tinggi itu disia-siakan.
\end{solution}

\begin{solution}{exo:b1:second-law-entropy:8}
Siklus: $\Delta S = 0 = Q/T_0 + S_{\mathrm{cipta}}$ dengan $S_{\mathrm{cipta}} \geq 0$,
sehingga $Q \leq 0$ dan $W = -Q \geq 0$: mesinnya tak dapat menyampaikan usaha.
Lautnya adalah satu termostat: kalornya tak dapat diubah menjadi usaha tanpa
benda yang lebih dingin untuk menampung entropinya.
\end{solution}

\begin{solution}{exo:b1:second-law-entropy:9}
Tiap gasnya memuai dari $V$ ke $2V$ pada $T$ tetap: $\Delta S = 2 \times R\ln2
= \qty{11.5}{J/K}$, tercipta (tak ada kalor yang dipertukarkan, kotaknya terkucil). Dua
separuh nitrogen: menyingkirkan dindingnya tidak mengubah apa pun yang bersifat makroskopik ---
$\Delta S = 0$ (paradoks Gibbs: molekul yang identik tidak ``tercampur'').
\end{solution}

\begin{solution}{exo:b1:second-law-entropy:10}
$Q = \qty{3.6e5}{J}$ masuk ke ruangan bersuhu \qty{293}{K}: $S = \qty{1230}{J/K}$
tercipta --- di dalam filamennya (usaha listrik yang berubah menjadi kalor) dan dalam
aliran kalor itu dari filamen yang panas ke ruangan yang sejuk; usaha
listriknya sendiri tak membawa entropi.
\end{solution}

\begin{solution}{exo:b1:second-law-entropy:11}
$\Delta S_{\text{benda}} = C\ln(T_0/T_1)$ pada tiap kasusnya. (a) yang dipertukarkan
$C(T_0 - T_1)/T_0$: yang tercipta $C[T_1/T_0 - 1 - \ln(T_1/T_0)]$. (b) dua tahap
yang masing-masing bernisbah $r = \sqrt{T_1/T_0}$: yang tercipta $2C[r - 1 - \ln r]$, lebih kecil
(untuk $T_1/T_0 = 2$: $0.193C$ berbanding $0.307C$). (c) tahap yang tak berhingga banyaknya:
masing-masing menciptakan $C[(1 + \epsilon) - 1 - \ln(1 + \epsilon)] \approx C\epsilon^2/2$,
yang berjumlah nol ketika $\epsilon \to 0$: pendinginan reversibel menuntut termostat
pada setiap suhu antaranya --- kalor tak boleh menyeberangi selisih yang berhingga.
\end{solution}

\begin{solution}{exo:b1:second-law-entropy:12}
$2^{-10} = 10^{-3}$; $2^{-100} = \num{8e-31}$; $2^{-10^{22}} = 10^{-3\times10^{21}}$.
$k_B\ln2^N = Nk_B\ln2 = nR\ln2$. ``Mustahil'' secara ketat berarti peluangnya
nol; di sini ia $10^{-3\times10^{21}}$ --- bukan nol, tetapi bilangan yang begitu kecil sehingga
tak ada bedanya yang pernah dapat diamati: ketakreversibelan adalah statistika pada
skala $10^{23}$.
\end{solution}

\begin{solution}{pb:b1:second-law-entropy:1}
\textbf{1.} $Q = mc(T_0 - T_1) = 0.25 \times 4180 \times (-70) = \qty{-73}{kJ}$
bagi kopinya: ruangannya memperoleh \qty{73}{kJ}, $\Delta S_{\text{ruangan}} = 73150/293
= \qty{+250}{J/K}$.

\textbf{2.} $mc\ln(T_0/T_1) = 1045\ln(293/363) = \qty{-224}{J/K}$.

\textbf{3.} $S_{\mathrm{cipta}} = -224 - (-250) = \qty{+26}{J/K} > 0$.

\textbf{4.} $\Delta S = 1045\ln(330/363) = \qty{-99.6}{J/K}$; $Q = 1045 \times
(-33) = \qty{-34.5}{kJ}$, $S_{\mathrm{tukar}} = \qty{-117.7}{J/K}$: tercipta
\qty{18}{J/K} sejauh ini --- sebagian besar ketakreversibelannya terjadi selagi
kopinya paling panas.

\textbf{5.} $\Delta S - Q/T_0 = mc[\ln(T_0/T_1) - (T_0 - T_1)/T_0] = mc[x - 1 -
\ln x]$ dengan $x = T_1/T_0$; $f'(x) = 1 - 1/x$ lenyap di $x = 1$ tempat
$f = 0$, dan $f'' = 1/x^2 > 0$: minimumnya $0$.

\textbf{6.} Tidak: keadaan awal dan akhirnya sama, kalornya
sama, ruangannya termostat yang sama --- entropi yang tercipta ditetapkan oleh
keadaannya, bukan oleh kecepatannya.

\textbf{7.} $T_f = \qty{350}{K}$; yang tercipta $4180\ln(350^2/120000) = \qty{86}{J/K}$.

\textbf{8.} Reversibel: entropi totalnya kekal: $mc\ln(T_f'/T_1) + mc\ln(T_f'/
T_2) = 0$, sehingga $T_f'^2 = T_1T_2$.

\textbf{9.} $T_f' = \sqrt{120000} = \qty{346.4}{K}$; kedua airnya kehilangan $mc[(T_1 -
T_f') + (T_2 - T_f')] = 4180 \times (53.6 - 46.4) = \qty{30}{kJ}$, yang disampaikan sebagai
usaha.

\textbf{10.} Yang panas: $4180\ln(346.4/400) = \qty{-601}{J/K}$; yang dingin: $4180\ln(346.4/
300) = \qty{+601}{J/K}$: jumlahnya nol, seperti yang dituntut kereversibelannya.

\textbf{11.} $T_0S_{\mathrm{cipta}} = 300 \times 86 = \qty{26}{kJ}$, dekat dengan
\qty{30}{kJ} yang dapat dipulihkan: entropi yang tercipta, dikali suhu
acuannya, adalah usaha yang dibuang oleh \omterm{def:b1:point-kinematics:frame}{lintasan} tak reversibelnya.

\textbf{12.} Ia akan menuntut sebuah mesin yang bekerja di antara dua air
suam-suam kuku yang suhunya terus berubah --- beberapa puluh kilojoule bagi
sebuah mesin besar: tak sepadan, tetapi hukum kedua mengatakan usahanya
ada di sana.

\textbf{13.} Yang diterima dari bagian dalamnya $Q_c/T_c = 40/278 = \qty{0.144}{W/K}$;
yang diberikan ke ruangannya $Q_0/T_0 = (40 + W)/293$; neracanya sepanjang satu siklus:
$0 = 0.144 - (40 + W)/293 + S_c$.

\textbf{14.} $S_c \geq 0$: $(40 + W)/293 \geq 0.144$, $W \geq 42.2 - 40 =
\qty{2.2}{W}$; COP $\leq 40/2.2 = 18$ ($= T_c/(T_0 - T_c)$).

\textbf{15.} $\mathrm{COP}_{\max} = T_c/(T_0 - T_c)$: bagi sebuah pembeku $255/38 = 6.7$
--- makin dingin kotaknya, makin mahal tiap joule yang dikeluarkan.

\textbf{16.} Dengan $W = 15$: $S_c = (55/293) - 0.144 = \qty{0.044}{W/K}$ tercipta
tiap sekon.

\textbf{17.} $W \geq 200 \times 15/278 = \qty{10.8}{W}$; mesinnya membuang
\qty{200}{W} ditambah usahanya sendiri ke dalam ruangannya sementara bagian dalamnya tinggal
dingin: perolehan total ruangannya adalah $W$ --- lemari es yang terbuka adalah sebuah pemanas.

\textbf{18.} $W = 0$ akan memberikan $S_c = Q_c(1/T_0 - 1/T_c) < 0$: terlarang ---
itulah \omterm{prop:b1:second-law-entropy:clausiuskelvin}{pernyataan Clausius}.

\textbf{19.} Sebelumnya: $1$ (semuanya di kiri); sesudahnya: $2^N$; $\Delta S = k_B\ln2^N =
Nk_B\ln2$.

\textbf{20.} $Nk_B = nR$: $\Delta S = nR\ln2$, yaitu hasil \omterm{rem:b1:first-law:irreversible}{ekspansi Joule}.

\textbf{21.} Satu bit tiap molekul: $N$ bit; dan $\Delta S/(k_B\ln2) = N$ ---
entropi mencacah informasi yang kita perlukan untuk memakukan molekulnya
di tempatnya.

\textbf{22.} $2^{-50} = \num{9e-16}$; bagi satu \omterm{def:b1:kinetic-theory:scales}{mol} $2^{-6\times10^{23}} =
10^{-1.8\times10^{23}}$.

\textbf{23.} Tak ada kalor ($Q = 0$): yang dipertukarkan $0$, yang tercipta $nR\ln2$ ---
kasus yang paling murni: tak ada usaha, tak ada kalor, hanya penyebaran gasnya.

\textbf{24.} $W = nRT\ln2$ diterima, $Q = -nRT\ln2$ diberikan ke ruangannya,
$S_{\mathrm{tukar}} = -nR\ln2$ (entropi gasnya kembali ke nilai
awalnya, secara reversibel): membatalkan ekspansi bebasnya menuntut usaha $nRT\ln2$
dan membuang $nR\ln2$ entropi ke dalam ruangannya --- semestanya menyimpan
$nR\ln2$ yang tercipta itu untuk selamanya.

\textbf{25.} Entropi tidak kekal pada apa pun yang nyata --- hanya pada
limit reversibel yang diidealkan; hukumnya melarang setiap proses yang akan
memusnahkannya (kalor yang mendaki, usaha dari satu termostat, penguraian campuran secara cuma-cuma);
dan tiap bit yang tercipta adalah usaha yang hilang, sebesar $T_0S_{\mathrm{cipta}}$.
\end{solution}
